\documentclass{article}

\usepackage[utf8]{inputenc}
\usepackage[T1]{fontenc}
\usepackage{helvet}
\usepackage{geometry}
\usepackage{xcolor}
\usepackage{eso-pic}
\usepackage{fancyhdr}
\usepackage{parskip}
\usepackage{microtype}
\usepackage[english, ukrainian]{babel}
\usepackage{url}
\usepackage{amsmath}
\usepackage{amssymb}
\usepackage{amsthm}
\usepackage{longtable}
\usepackage{booktabs}
\usepackage{hyperref}
\usepackage{titlesec}
\usepackage{tocloft}

\renewcommand{\cftsecfont}{\small\bfseries}
\renewcommand{\cftsubsecfont}{\footnotesize}
\renewcommand{\cftsubsubsecfont}{\scriptsize}
\renewcommand{\cftsecpagefont}{\small\bfseries}
\renewcommand{\cftsubsecpagefont}{\footnotesize}
\renewcommand{\cftsubsubsecpagefont}{\scriptsize}

\definecolor{nato}{HTML}{293757}
\definecolor{footergray}{gray}{0.65}

\geometry{a4paper,left=3cm,right=3cm,top=3cm,bottom=3cm}
\renewcommand{\familydefault}{\sfdefault}
\setcounter{secnumdepth}{3}

\titleformat{\section}
  {\color{nato}\Huge\bfseries\raggedright}{\thesection.}{1em}{}
\titlespacing*{\section}{0pt}{30pt}{12pt}

\titleformat{\subsection}
  {\color{nato}\LARGE\bfseries\raggedright}{\thesubsection.}{1em}{}
\titlespacing*{\subsection}{0pt}{20pt}{8pt}

\titleformat{\subsubsection}
  {\color{nato}\Large\bfseries\raggedright}{\thesubsubsection.}{1em}{}
\titlespacing*{\subsubsection}{0pt}{10pt}{6pt}

\fancyhf{}
\fancyhead[R]{\small\color{footergray} 2 червня 2026}
\fancyfoot[L]{\small\color{footergray} Комплексна система захисту інформації. Завдання безпеки для військової пошти (20).}
\fancyfoot[R]{\small\color{footergray}\thepage}

\newcommand\HUGE[1]{\fontsize{40}{40}\selectfont #1}

\renewcommand{\headrulewidth}{0pt}
\renewcommand{\footrulewidth}{0.4pt}
\renewcommand{\footrule}{\color{footergray}\hrule width \textwidth height 0.4pt}

\begin{document}

\pagestyle{fancy}

\begin{titlepage}
  \AddToShipoutPictureBG*{\AtPageLowerLeft{\color{nato}\rule{\paperwidth}{\paperheight}}}
  \color{white}\thispagestyle{empty}\vspace*{4cm}\centering
  {\Huge  \textbf{Комплексна система захисту інформації} \par} \vspace{2cm}
  {\HUGE  \textbf{Завдання безпеки для військової пошти (20)} \par} \vspace{2.5cm}
  \vfill
  {\large 2026 \copyright\ Критпографічні Телесистеми \par}
\end{titlepage}

\newpage

\begin{center}
  {\Huge\color{nato}\bfseries Зміст\par}
\end{center}
\vspace{40pt}
\tableofcontents

\begin{abstract}
Завдання з безпеки для забезпечення конфіденційності, цілісності та гарантованої доставки повідомлень (MHS X.420) з підтримкою грифування.
\end{abstract}

\section{AC}
Клас заходів захисту AC --- УПРАВЛІННЯ ДОСТУПОМ

Цей клас зосереджується на обмеженні доступу до інформаційних систем, ресурсів та функцій лише для авторизованих користувачів, програм і пристроїв.

\textbf{Перелік заходів захисту:} Забезпечення доступу (AC-3); Управління інформаційними потоками (AC-4); Управління паралельною сесією (AC-10); Атрибути безпеки та приватності (AC-16); Id-spe-ac-16-5; Диспетчер доступу (AC-25).

\subsection{ЗАБЕЗПЕЧЕННЯ ДОСТУПУ (AC-3)}
Застосовувати затверджені повноваження для логічного доступу до інформації та ресурсів системи відповідно до чинної політики (правил) управління доступом.

\vspace{10pt}
{\footnotesize\bfseries
No: 1\\
Name: ac\_3\_01\\
Type: list\\
Default: [\textquotedbl{}default\_deny\_rule\textquotedbl{}, \textquotedbl{}abac\_rule\_1\textquotedbl{}]
}

{\footnotesize Затверджені повноваження на логічний доступ до інформації та ресурсів системи виконуються відповідно до чинних політик(правил) управління доступом}




\subsection{УПРАВЛІННЯ ІНФОРМАЦІЙНИМИ ПОТОКАМИ (AC-4)}
Застосувати затверджені повноваження для управління потоком інформації всередині системи та між пов'язаними системами на основі [Призначення: визначеними організацією політиками управління інформаційним потоком].

\vspace{10pt}
{\footnotesize\bfseries
No: 1\\
Name: ac\_4\_01\\
Type: list\\
Default: [\textquotedbl{}default\_deny\_rule\textquotedbl{}, \textquotedbl{}abac\_rule\_1\textquotedbl{}]
}

{\footnotesize Затверджені повноваження застосовуються для контролю потоку інформації всередині системи та між підключеними системами на основі політики управління інформаційними потоками}


{\footnotesize\bfseries
No: 2\\
Name: ac\_4\_odp\\
Type: list\\
Default: [\textquotedbl{}default\_deny\_rule\textquotedbl{}, \textquotedbl{}abac\_rule\_1\textquotedbl{}]
}

{\footnotesize Визначено політики управління інформаційними потоками всередині системи та між підключеними системами}




\subsection{УПРАВЛІННЯ ПАРАЛЕЛЬНОЮ СЕСІЄЮ (AC-10)}
Обмежити кількість одночасних сеансів для кожного [Призначення: визначеного організацією облікового запису та/або типу облікового запису] до [Призначення: визначеної організацією кількості].

\vspace{10pt}
{\footnotesize\bfseries
No: 1\\
Name: ac\_10\_01\\
Type: integer\\
Default: 30
}

{\footnotesize Кількість одночасних сеансів для кожного облікового запису та/або типів облікових записів обмежена кількість}


{\footnotesize\bfseries
No: 2\\
Name: ac\_10\_odp\_02\\
Type: integer\\
Default: 30
}

{\footnotesize Визначено кількість одночасних сеансів, дозволених для кожного облікового запису та/або типу облікового запису}




\subsection{АТРИБУТИ БЕЗПЕКИ ТА ПРИВАТНОСТІ (AC-16)}
a. Визначити засоби для асоціювання (пов'язання) [Призначення: визначених організацією типів атрибутів безпеки та приватності], що приймають [Призначення: визначені організацією значення атрибутів безпеки та приватності] з інформацією, яка зберігається, обробляється та/або передається.\\ 
b. Пов'язані атрибути безпеки та приватності мають створюватися і зберігатися разом з інформацією.\\ 
c. Встановити дозволені [Призначення: визначені організацією атрибути безпеки] для [Призначення: систем, визначених організацією].\\ 
d. Визначити дозволені [Призначення: визначені організацією значення або діапазони] для кожного з встановлених атрибутів безпеки та приватності.

\vspace{10pt}
{\footnotesize\bfseries
No: 1\\
Name: ac\_16\_c\_02\\
Type: list\\
Default: [\textquotedbl{}default\_deny\_rule\textquotedbl{}, \textquotedbl{}abac\_rule\_1\textquotedbl{}]
}

{\footnotesize Визначено типи атрибутів безпеки, які мають бути пов'язані зі значеннями атрибутів безпеки для інформації, що зберігається, обробляється та/або передається; визначено значення атрибутів конфіденційності для типів атрибутів конфіденційності; визначено системи, для яких мають бути встановлені дозволені атрибути безпеки; визначено системи, для яких мають бути встановлені дозволені атрибути конфіденційності; визначено атрибути безпеки, визначені як частина AC-16(a), які дозволені для систем; визначено атрибути конфіденційності, визначені як частина AC-16(a), які дозволені для систем; визначено значення атрибутів або діапазони для встановлених атрибутів; визначено частоту, з якою слід переглядати атрибути безпеки на предмет відповідності; визначено частоту, з якою слід переглядати атрибути конфіденційності на предмет відповідності; надано засоби для асоціювання типів атрибутів безпеки зі значеннями атрибутів безпеки для інформації, що зберігається, обробляється та/або передається надано засоби для асоціювання типів атрибутів конфіденційності зі значеннями атрибутів конфіденційності для інформації, що зберігається, обробляється та/або передається виникають зв'язки з атрибутами; зв'язки атрибутів зберігаються разом з інформацією; на основі атрибутів, визначених у AC-16\_ODP[01] для систем, встановлюються наступні дозволені атрибути безпеки: атрибути безпеки; на основі атрибутів, визначених у AC-16\_ODP[02] для систем, встановлюються наступні дозволені атрибути приватності: атрибути конфіденційності ; AC-16(d)}


{\footnotesize\bfseries
No: 2\\
Name: ac\_16\_f\_02\\
Type: list\\
Default: [\textquotedbl{}default\_deny\_rule\textquotedbl{}, \textquotedbl{}abac\_rule\_1\textquotedbl{}]
}

{\footnotesize Визначено наступні допустимі значення або діапазони атрибутів для кожного з встановлених атрибутів: значення або діапазони атрибутів; проводиться аудит змін до атрибутів; атрибути безпеки перевіряються на відповідність частота; атрибути конфіденційності перевіряються на відповідність частота}


{\footnotesize\bfseries
No: 3\\
Name: ac\_16\_odp\_02\\
Type: string\\
Default: nil
}

{\footnotesize Визначено типи атрибутів конфіденційності, які мають бути пов'язані зі значеннями атрибутів конфіденційності для інформації, що зберігається, обробляється та/або передається}


{\footnotesize\bfseries
No: 4\\
Name: ac\_16\_odp\_03\\
Type: string\\
Default: nil
}

{\footnotesize Визначено значення атрибутів безпеки для типів атрибутів безпеки}




\subsubsection{id-spe-ac-16-5}
Відображати атрибути безпеки та приватності в зручній для людини формі для кожного об'єкту, який система передає на пристрої виведення, щоб ідентифікувати [Призначення: визначені організацією спеціальні інструкції щодо поширення, обробки чи наступного розподілу інформації], використовуючи [Призначення: визначену організацією ідентифікацію, у зручній для людини формі про стандартні угоди про присвоєння імен].

\vspace{10pt}
{\footnotesize\bfseries
No: 1\\
Name: ac\_16\_5\_02\\
Type: list\\
Default: [\textquotedbl{}default\_deny\_rule\textquotedbl{}, \textquotedbl{}abac\_rule\_1\textquotedbl{}]
}

{\footnotesize Для кожного об'єкта, який система передає на пристрої виведення, визначено спеціальні інструкції щодо поширення, обробки чи розподілу, які мають використовуватися; визначено стандартні угоди про ідентифікацію атрибутів безпеки та конфіденційності, які повинні відображатися в зручній для читання формі на кожному об'єкті, який система передає на пристрої виводу; атрибути безпеки відображаються у формі, зручній для читання людиною, на кожному об'єкті, який система передає на пристрої виводу для ідентифікації інструкцій з використанням угод про іменування; атрибути конфіденційності відображаються у формі, зручній для читання людиною, на кожному об'єкті, який система передає на пристрої виводу для ідентифікації інструкцій з використанням угод про іменування}




\subsection{ДИСПЕТЧЕР ДОСТУПУ (AC-25)}
Впровадити диспетчер доступу для [Призначення: визначеної організацією політики контролю доступу], який захищений від несанкціонованого доступу, завжди був доступний для виклику та досить компактний, щоб бути підданим аналізу й тестуванню, надійність якого може бути гарантована.

\vspace{10pt}
{\footnotesize\bfseries
No: 1\\
Name: ac\_25\_01\\
Type: list\\
Default: [\textquotedbl{}default\_deny\_rule\textquotedbl{}, \textquotedbl{}abac\_rule\_1\textquotedbl{}]
}

{\footnotesize Реалізовано диспетчер доступу для політики контролю доступу, який захищений від несанкціонованого доступу, завжди доступний для виклику та досить компактний, щоб бути підданим аналізу й тестуванню, надійність якого може бути гарантована}


{\footnotesize\bfseries
No: 2\\
Name: ac\_25\_odp\\
Type: list\\
Default: [\textquotedbl{}default\_deny\_rule\textquotedbl{}, \textquotedbl{}abac\_rule\_1\textquotedbl{}]
}

{\footnotesize Визначено політики контролю доступу, для яких реалізовано диспетчер доступу}




\section{AU}
Клас заходів захисту AU --- АУДИТ ТА ПІДЗВІТНІСТЬ

Цей клас забезпечує можливість відстеження дій у системі, генерування, захист та аналіз записів аудиту для виявлення порушень політики безпеки.

\textbf{Перелік заходів захисту:} Події аудиту (AU-2); Неспростовність (AU-10); Id-spe-au-10-5; Генерація даних аудиту (AU-12).

\subsection{ПОДІЇ АУДИТУ (AU-2)}
a. Визначити типи подій, які система може реєструвати для підтримки функції аудиту: [Призначення: типи подій, визначені організацією, які система здатна реєструвати];\\ 
b. Координувати функції аудиту безпеки з іншими організаційними підрозділами, які вимагають інформації, пов'язаної з аудитом, для посилення взаємної підтримки та допомоги у виборі типів подій, що перевіряються;\\ 
c. Визначити, які типи подій підлягають аудиту: [Призначення: визначені організацією події, що підлягають аудиту (підмножина подій, що підлягають аудиту, визначених в AU-2 a.), а також частота (або ситуація, що вимагає) проведення аудиту для кожної ідентифікованої події]\\ 
d. Обґрунтувати, чому типи подій, що перевіряються, вважаються достатніми для підтримки розслідувань інцидентів (постфактум), пов'язаних з безпекою та приватністю;\\ 
e. Перегляньте й оновіть типи подій, вибрані для журналювання [Призначення: частота, визначена організацією].

\vspace{10pt}
{\footnotesize\bfseries
No: 1\\
Name: au\_2\_a\\
Type: string\\
Default: nil
}

{\footnotesize Типи подій, які система здатна реєструвати, визначено для підтримки функції аудиту}


{\footnotesize\bfseries
No: 2\\
Name: au\_2\_b\\
Type: string\\
Default: nil
}

{\footnotesize Функція аудиту безпеки координується з іншими підрозділами організації, які вимагають інформації, пов'язаної з аудитом, длдля посилення взаємної підтримки та допомоги у виборі типів подій, що перевіряються}


{\footnotesize\bfseries
No: 3\\
Name: au\_2\_c\_01\\
Type: string\\
Default: nil
}

{\footnotesize Типи подій (підмножина 02\_ODP[01]) визначаються для реєстрації у системі; AU-}


{\footnotesize\bfseries
No: 4\\
Name: au\_2\_c\_02\\
Type: string\\
Default: \textquotedbl{}щорічно\textquotedbl{}
}

{\footnotesize Зазначені типи подій реєструються 02\_ODP[03] частота або ситуація>; <AU-}


{\footnotesize\bfseries
No: 5\\
Name: au\_2\_d\\
Type: string\\
Default: nil
}

{\footnotesize Надається обґрунтування, чому типи подій, що перевіряються, вважаються достатніми для підтримки розслідувань інцидентів (постфактум), пов'язаних з безпекою та конфіденційністю}


{\footnotesize\bfseries
No: 6\\
Name: au\_2\_e\\
Type: string\\
Default: \textquotedbl{}щорічно\textquotedbl{}
}

{\footnotesize Переглядаються та оновлюються типи подій, вибрані для реєстрації, частота. у системі}


{\footnotesize\bfseries
No: 7\\
Name: au\_2\_odp\_01\\
Type: string\\
Default: nil
}

{\footnotesize Визначено типи подій, які система може реєструвати для підтримки функції аудиту}


{\footnotesize\bfseries
No: 8\\
Name: au\_2\_odp\_02\\
Type: string\\
Default: nil
}

{\footnotesize Визначено типи подій (підмножина AU-02\_ODP[01]) що підлягають аудиту у системі}


{\footnotesize\bfseries
No: 9\\
Name: au\_2\_odp\_03\\
Type: list\\
Default: [\textquotedbl{}login\textquotedbl{}, \textquotedbl{}logout\textquotedbl{}, \textquotedbl{}failed\_attempt\textquotedbl{}]
}

{\footnotesize Визначено частоту або ситуацію, що вимагає проведення аудиту для кожної ідентифікованої події}


{\footnotesize\bfseries
No: 10\\
Name: au\_2\_odp\_04\\
Type: string\\
Default: \textquotedbl{}щорічно\textquotedbl{}
}

{\footnotesize Частота перегляду та оновлення типів подій, обраних для журналювання}




\subsection{НЕСПРОСТОВНІСТЬ (AU-10)}
Надавайте неспростовні докази того, що особа (або процес, який діє від імені особи) виконала [Призначення: дії, визначені організацією, на які поширюється принцип неспростовності].

\vspace{10pt}
{\footnotesize\bfseries
No: 1\\
Name: au\_10\_01\\
Type: list\\
Default: [\textquotedbl{}admin\textquotedbl{}, \textquotedbl{}security\_officer\textquotedbl{}]
}

{\footnotesize Надаються неспростовні докази того, що особа (або процес, що діє від імені особи) виконала дії}


{\footnotesize\bfseries
No: 2\\
Name: au\_10\_odp\\
Type: list\\
Default: [\textquotedbl{}login\textquotedbl{}, \textquotedbl{}logout\textquotedbl{}, \textquotedbl{}failed\_attempt\textquotedbl{}]
}

{\footnotesize Визначено дії, на які поширюється принцип неспростовності}




\subsubsection{id-spe-au-10-5}


\vspace{10pt}
{\footnotesize\bfseries
No: 1\\
Name: au\_10\_5\_01\\
Type: string\\
Default: nil
}

{\footnotesize НЕСПРОСТОВНІСТЬ - ЦИФРОВІ ПІДПИСИ [Вилучено: Включено до SI-07]}




\subsection{ГЕНЕРАЦІЯ ДАНИХ АУДИТУ (AU-12)}
a. Забезпечити генерацію даних аудиту для типів подій, що перевіряються в AU-2а в [Призначення: визначених організацією компонентах системи].\\ 
b. Дозволити [Призначення: визначеному організацією персоналу або посадам] вибирати, які типи подій, що перевіряються, повинні перевірятися окремими компонентами системи;\\ 
c. Генерувати записи аудиту для типів подій, визначених в AU-2с. з вмістом згідно з AU-3.

\vspace{10pt}
{\footnotesize\bfseries
No: 1\\
Name: au\_12\_a\\
Type: string\\
Default: nil
}

{\footnotesize Можливість генерації записів аудиту для типів подій, які система здатна перевіряти (визначених у AU-02\_ODP[01]), забезпечується компонентами системи}


{\footnotesize\bfseries
No: 2\\
Name: au\_12\_b\\
Type: list\\
Default: [\textquotedbl{}admin\textquotedbl{}, \textquotedbl{}security\_officer\textquotedbl{}]
}

{\footnotesize Персонал або ролі може/можуть вибирати типи подій, які будуть реєструватися певними компонентами системи; AU-12(c) згенеровано записи аудиту для типів подій, визначених у AU02\_ODP[02], які включають вміст записів аудиту, визначений у AU03}


{\footnotesize\bfseries
No: 3\\
Name: au\_12\_odp\_01\\
Type: string\\
Default: nil
}

{\footnotesize Визначено компоненти системи, які забезпечують можливість генерації записів аудиту для типів подій (визначених у AU02\_ODP[02])}


{\footnotesize\bfseries
No: 4\\
Name: au\_12\_odp\_02\\
Type: list\\
Default: [\textquotedbl{}admin\textquotedbl{}, \textquotedbl{}security\_officer\textquotedbl{}]
}

{\footnotesize Визначено персонал або ролі, яким дозволено обирати типи подій, що мають реєструватися певними компонентами системи}




\section{CM}
Клас заходів захисту CM --- УПРАВЛІННЯ КОНФІГУРАЦІЄЮ

Цей клас гарантує створення та підтримання базових налаштувань системи, а також строгий контроль за змінами в апаратному та програмному забезпеченні.

\textbf{Перелік заходів захисту:} Обмеження доступу до зміни (CM-5).

\subsection{ОБМЕЖЕННЯ ДОСТУПУ ДО ЗМІНИ (CM-5)}
Визначити, задокументувати, затвердити та забезпечити застосування фізичних і логічних обмежень доступу, пов'язаних зі змінами в системі.

\vspace{10pt}
{\footnotesize\bfseries
No: 1\\
Name: cm\_5\_01\\
Type: string\\
Default: nil
}

{\footnotesize Визначені та задокументовані фізичні обмеження доступу, пов'язані зі змінами в системі}


{\footnotesize\bfseries
No: 2\\
Name: cm\_5\_02\\
Type: string\\
Default: nil
}

{\footnotesize Затверджені фізичні обмеження доступу, пов'язані зі змінами в системі}


{\footnotesize\bfseries
No: 3\\
Name: cm\_5\_03\\
Type: string\\
Default: nil
}

{\footnotesize Застосовуються фізичні обмеження доступу, пов'язані зі змінами в системі}


{\footnotesize\bfseries
No: 4\\
Name: cm\_5\_04\\
Type: string\\
Default: nil
}

{\footnotesize Визначені та задокументовані логічні обмеження доступу, пов'язані зі змінами в системі}


{\footnotesize\bfseries
No: 5\\
Name: cm\_5\_05\\
Type: string\\
Default: nil
}

{\footnotesize Затверджені логічні обмеження доступу, пов'язані зі змінами в системі}


{\footnotesize\bfseries
No: 6\\
Name: cm\_5\_06\\
Type: string\\
Default: nil
}

{\footnotesize Застосовуються логічні обмеження доступу, пов'язані зі змінами в системі}




\section{IA}
Клас заходів захисту IA --- ІДЕНТИФІКАЦІЯ ТА АВТЕНТИФІКАЦІЯ

Цей клас відповідає за однозначне розпізнавання користувачів або пристроїв та підтвердження їхньої справжності перед наданням доступу до системи.

\textbf{Перелік заходів захисту:} Id-spe-ia-2-12; Id-spe-ia-5-11.

\subsubsection{id-spe-ia-2-12}


\vspace{10pt}
{\footnotesize\bfseries
No: 1\\
Name: ia\_2\_12\_01\\
Type: list\\
Default: [\textquotedbl{}admin\textquotedbl{}, \textquotedbl{}security\_officer\textquotedbl{}]
}

{\footnotesize Приймаються та електронним шляхом підтверджуються повноваження облікових даних особистої ідентифікації}




\subsubsection{id-spe-ia-5-11}


\vspace{10pt}
{\footnotesize\bfseries
No: 1\\
Name: ia\_5\_11\_01\\
Type: string\\
Default: nil
}

{\footnotesize УПРАВЛІННЯ АВТЕНТИФІКАТОРОМ - АВТЕНТИФІКАЦІЯ НА ОСНОВІ АПАРАТНИХ ТОКЕНІВ [Вилучено: Включено до ІА-02(01), ІА-02(02)]}




\section{MP}
Клас заходів захисту MP --- ЗАХИСТ НОСІЇВ ІНФОРМАЦІЇ

Цей клас спрямований на безпечне зберігання, транспортування, використання та знищення як цифрових, так і паперових носіїв інформації.

\textbf{Перелік заходів захисту:} Id-spe-mp-6; Id-spe-mp-8-4.

\subsection{id-spe-mp-6}
a. Очищувати [Призначення: визначені організацією системні носії] перед утилізацією, випуском за межі організаційного контролю, або перед повторним використанням [Призначення: методами та процедурами очищення, визначеними організацією].\\ 
b. Використовувати механізми очищення зі стійкістю та цілісністю, що відповідає категорії безпеки або рівню секретності інформації.

\vspace{10pt}
{\footnotesize\bfseries
No: 1\\
Name: mp\_6\_a\_01\\
Type: string\\
Default: nil
}

{\footnotesize Носії інформації системи перед утилізацією піддаються очищенню за допомогою методів та процедур очищення}


{\footnotesize\bfseries
No: 2\\
Name: mp\_6\_a\_02\\
Type: string\\
Default: nil
}

{\footnotesize Носії інформації системи очищаються за допомогою методів та процедур очищення перед випуском за межі контрольованої зони}


{\footnotesize\bfseries
No: 3\\
Name: mp\_6\_a\_03\\
Type: string\\
Default: nil
}

{\footnotesize Носії інформації системи очищуються за допомогою методів і процедур очищення перед повторним використанням}


{\footnotesize\bfseries
No: 4\\
Name: mp\_6\_b\\
Type: string\\
Default: \textquotedbl{}автоматизований засіб моніторингу\textquotedbl{}
}

{\footnotesize Застосовуються механізми очищення, надійність і цілісність яких відповідає категорії безпеки або рівню секретності інформації}


{\footnotesize\bfseries
No: 5\\
Name: mp\_6\_odp\_01\\
Type: string\\
Default: nil
}

{\footnotesize Визначено носії інформації системи, які підлягають очищенню перед утилізацією}


{\footnotesize\bfseries
No: 6\\
Name: mp\_6\_odp\_02\\
Type: string\\
Default: nil
}

{\footnotesize Визначено носії інформації системи, які підлягають очищенню перед випуском за межі контрольованої зони}


{\footnotesize\bfseries
No: 7\\
Name: mp\_6\_odp\_03\\
Type: string\\
Default: nil
}

{\footnotesize Визначені носії інформації системи, що підлягають очищенню перед повторним використанням}


{\footnotesize\bfseries
No: 8\\
Name: mp\_6\_odp\_04\\
Type: string\\
Default: nil
}

{\footnotesize Визначено методи та процедури очищення, які слід використовувати для очищення перед утилізацією}


{\footnotesize\bfseries
No: 9\\
Name: mp\_6\_odp\_05\\
Type: string\\
Default: nil
}

{\footnotesize Визначено методи та процедури очищення, які слід використовувати для очищення перед випуском за межі контрольованої зони}


{\footnotesize\bfseries
No: 10\\
Name: mp\_6\_odp\_06\\
Type: string\\
Default: nil
}

{\footnotesize Визначено методи та процедури очищення, які слід використовувати для очищення перед повторним використанням}




\subsubsection{id-spe-mp-8-4}


\vspace{10pt}
{\footnotesize\bfseries
No: 1\\
Name: mp\_8\_4\_01\\
Type: string\\
Default: nil
}

{\footnotesize Ідентифіковано носії інформації, що містять інформацію з обмеженим доступом}


{\footnotesize\bfseries
No: 2\\
Name: mp\_8\_4\_02\\
Type: list\\
Default: [\textquotedbl{}admin\textquotedbl{}, \textquotedbl{}security\_officer\textquotedbl{}]
}

{\footnotesize Носії інформації, що містять інформацію з обмеженим доступом, знижуються в класі перед передачею особам, які не мають необхідних дозволів на доступ}




\section{SC}
Клас заходів захисту SC --- ЗАХИСТ ІНФОРМАЦІЙНОЇ СИСТЕМИ ТА

Цей клас охоплює механізми захисту інформації під час її передачі мережами зв'язку, криптографічний захист та ізоляцію критичних компонентів.

\textbf{Перелік заходів захисту:} Id-spe-sc-8-1; Id-spe-sc-8-2; Встановлення ключами (SC-12); Id-spe-sc-17; Id-spe-sc-28-1.

\subsubsection{id-spe-sc-8-1}


\vspace{10pt}
{\footnotesize\bfseries
No: 1\\
Name: sc\_8\_1\_odp\\
Type: string\\
Default: nil
}

{\footnotesize Вибрано одне або декілька з наступних ЗНАЧЕНЬ ПАРАМЕТРІВ: \{запобігти несанкціонованому розголошенню інформації; виявити зміни в інформації\}}




\subsubsection{id-spe-sc-8-2}


\vspace{10pt}
{\footnotesize\bfseries
No: 1\\
Name: sc\_8\_2\_01\\
Type: string\\
Default: nil
}

{\footnotesize КОНФІДЕНЦІЙНІСТЬ ТА ЦІЛІСНІСТЬ ПЕРЕДАЧІ - ПОПЕРЕДНЯ І ПОСТОБРОБКА МЕТА ОЦІНКИ: Визначити, чи:}


{\footnotesize\bfseries
No: 2\\
Name: sc\_8\_2\_odp\\
Type: string\\
Default: nil
}

{\footnotesize Вибрано одне або декілька з наступних ЗНАЧЕНЬ ПАРАМЕТРІВ: \{конфіденційність; цілісність\}}




\subsection{ВСТАНОВЛЕННЯ КЛЮЧАМИ (SC-12)}
Встановити та управляти криптографічними ключами для криптографічних засобів, які використовуються в системі відповідно до [Призначення: визначені організацією вимоги до генерації, поширення, зберігання, доступу та знищення ключів].

\vspace{10pt}
{\footnotesize\bfseries
No: 1\\
Name: sc\_12\_01\\
Type: string\\
Default: \textquotedbl{}AES-256-GCM\textquotedbl{}
}

{\footnotesize Встановлюються криптографічні ключі, коли в системі використовується криптографія відповідно до < SC-12\_ODP вимог >}


{\footnotesize\bfseries
No: 2\\
Name: sc\_12\_02\\
Type: string\\
Default: \textquotedbl{}AES-256-GCM\textquotedbl{}
}

{\footnotesize Здійснюється управління криптографічними ключами, коли в системі використовується криптографія, відповідно до вимог }


{\footnotesize\bfseries
No: 3\\
Name: sc\_12\_odp\\
Type: string\\
Default: nil
}

{\footnotesize Визначені вимоги до генерації, розповсюдження, зберігання, доступу та знищення ключів}




\subsection{id-spe-sc-17}
a. Випускати сертифікати відкритого ключа відповідно до [Призначення: визначеної організацією політики сертифікації];\\ 
b. Отримувати сертифікати відкритого ключа від затвердженого постачальника послуг.

\vspace{10pt}
\textit{Немає параметрів для цього контролю.}
\vspace{10pt}


\subsubsection{id-spe-sc-28-1}


\vspace{10pt}
{\footnotesize\bfseries
No: 1\\
Name: sc\_28\_1\_01\\
Type: string\\
Default: \textquotedbl{}AES-256-GCM\textquotedbl{}
}

{\footnotesize Реалізовані криптографічні механізми для запобігання несанкціонованому розкриттю інформації, що знаходиться в стані спокою на системних компонентах або носіях}


{\footnotesize\bfseries
No: 2\\
Name: sc\_28\_1\_02\\
Type: string\\
Default: \textquotedbl{}AES-256-GCM\textquotedbl{}
}

{\footnotesize Реалізовані криптографічні механізми для запобігання несанкціонованій модифікації інформації, що знаходиться в стані спокою на системних компонентах або носіях}




\end{document}
